\documentclass[11pt,a4paper]{article}

\usepackage{amsmath,amssymb,amsthm}
\usepackage{graphicx}
\usepackage{hyperref}
\usepackage{algorithm}
\usepackage{algpseudocode}
\usepackage{booktabs}
\usepackage{geometry}
\usepackage{authblk}

\geometry{margin=1in}

\theoremstyle{definition}
\newtheorem{definition}{Definition}
\newtheorem{theorem}{Theorem}
\newtheorem{lemma}{Lemma}
\newtheorem{proposition}{Proposition}
\newtheorem{corollary}{Corollary}

\title{Decentralized Lending Protocols: \\
Collateralization, Liquidation, and Interest Rate Dynamics}

\author{Zach Kelling}
\affil{Hanzo Industries \\ Lux Industries \\ Zoo Labs Foundation}
\date{2022}

\begin{document}

\maketitle

\begin{abstract}
We present a comprehensive mathematical framework for decentralized lending protocols on the Lux blockchain. Our analysis covers three fundamental components: collateralization mechanics that ensure protocol solvency, liquidation mechanisms that maintain health of under-collateralized positions, and dynamic interest rate models that balance supply and demand for liquidity. We derive optimal parameters for each component, prove safety properties, and provide implementation guidance for high-throughput blockchain environments. The framework enables permissionless, over-collateralized lending with provable safety guarantees.
\end{abstract}

\section{Introduction}

Decentralized lending protocols enable permissionless borrowing and lending of digital assets. Unlike traditional finance, these protocols operate without trusted intermediaries, relying instead on smart contracts and economic incentives.

The core challenge is maintaining solvency: ensuring that outstanding loans are always backed by sufficient collateral, even under adverse market conditions. This requires:

\begin{enumerate}
\item \textbf{Collateralization}: Borrowers must deposit assets exceeding their loan value
\item \textbf{Liquidation}: Under-collateralized positions must be resolved promptly
\item \textbf{Interest Rates}: Rates must adapt to balance supply and demand
\end{enumerate}

This paper provides rigorous mathematical foundations for each component.

\section{Collateralization Framework}

\subsection{Position Definition}

\begin{definition}[Lending Position]
A lending position is a tuple $(C, D, t)$ where:
\begin{itemize}
\item $C$: Collateral value in reference currency
\item $D$: Debt value in reference currency
\item $t$: Timestamp of last update
\end{itemize}
\end{definition}

\begin{definition}[Health Factor]
The health factor of a position is:
\begin{equation}
H = \frac{C \cdot \theta}{D}
\end{equation}
where $\theta \in (0, 1)$ is the liquidation threshold for the collateral asset.
\end{definition}

\begin{definition}[Loan-to-Value Ratio]
The loan-to-value (LTV) ratio is:
\begin{equation}
\text{LTV} = \frac{D}{C}
\end{equation}
\end{definition}

\subsection{Collateral Factors}

Different assets have different risk profiles, requiring asset-specific parameters:

\begin{definition}[Asset Parameters]
For each supported collateral asset $a$, we define:
\begin{itemize}
\item $\theta_a$: Liquidation threshold (e.g., 0.80 for 80\%)
\item $\lambda_a$: Maximum LTV for new borrows (e.g., 0.75)
\item $\beta_a$: Liquidation bonus (e.g., 0.05 for 5\%)
\end{itemize}
\end{definition}

\begin{proposition}[Parameter Constraints]
For protocol safety:
\begin{equation}
\lambda_a < \theta_a < 1 - \beta_a
\end{equation}
\end{proposition}

\begin{proof}
\begin{itemize}
\item $\lambda_a < \theta_a$: New borrows must have buffer before liquidation
\item $\theta_a < 1 - \beta_a$: Liquidators must profit after receiving bonus
\end{itemize}
\end{proof}

\subsection{Multi-Asset Positions}

\begin{definition}[Aggregate Health Factor]
For a position with multiple collateral types and debts:
\begin{equation}
H = \frac{\sum_i C_i \cdot \theta_i}{\sum_j D_j}
\end{equation}
\end{definition}

\begin{theorem}[Solvency Condition]
A position is solvent if and only if $H \geq 1$.
\end{theorem}

\section{Liquidation Mechanism}

\subsection{Liquidation Trigger}

\begin{definition}[Liquidatable Position]
A position is liquidatable when $H < 1$, equivalently:
\begin{equation}
\sum_j D_j > \sum_i C_i \cdot \theta_i
\end{equation}
\end{definition}

\subsection{Liquidation Process}

\begin{algorithm}
\caption{Liquidation Execution}
\begin{algorithmic}[1]
\Require Liquidatable position $(C, D)$, repay amount $R$, liquidation bonus $\beta$
\Ensure Updated position, liquidator profit
\State \textbf{require} $H < 1$
\State $R_{\max} \gets \min(D \cdot \phi, D)$ \Comment{$\phi$ is close factor, e.g., 0.5}
\State \textbf{require} $R \leq R_{\max}$
\State $\text{Collateral Seized} \gets R \cdot (1 + \beta) / P_C$
\State \textbf{require} $\text{Collateral Seized} \leq C$
\State $D' \gets D - R$
\State $C' \gets C - \text{Collateral Seized}$
\State \Return $(C', D')$
\end{algorithmic}
\end{algorithm}

\subsection{Close Factor}

\begin{definition}[Close Factor]
The close factor $\phi \in (0, 1]$ limits the maximum debt repayable in a single liquidation:
\begin{equation}
R_{\max} = D \cdot \phi
\end{equation}
\end{definition}

\begin{proposition}[Close Factor Rationale]
A close factor $\phi < 1$ (e.g., 0.5) provides:
\begin{enumerate}
\item Multiple liquidation opportunities, increasing competition
\item Protection against over-liquidation
\item Gas efficiency for partial liquidations
\end{enumerate}
\end{proposition}

\subsection{Liquidation Incentives}

\begin{theorem}[Liquidator Profitability]
A liquidator repaying $R$ receives collateral worth $R(1 + \beta)$. The profit is:
\begin{equation}
\Pi = R \cdot \beta - G
\end{equation}
where $G$ is gas cost in reference currency.
\end{theorem}

\begin{corollary}[Minimum Profitable Liquidation]
Liquidation is profitable when:
\begin{equation}
R > \frac{G}{\beta}
\end{equation}
\end{corollary}

\subsection{Bad Debt Handling}

\begin{definition}[Bad Debt]
Bad debt occurs when $C < D$ after liquidation:
\begin{equation}
\text{Bad Debt} = \max(D - C, 0)
\end{equation}
\end{definition}

\begin{proposition}[Bad Debt Prevention]
Bad debt is prevented when:
\begin{equation}
\theta < \frac{1}{1 + \beta}
\end{equation}
This ensures liquidation is always profitable before insolvency.
\end{proposition}

\section{Interest Rate Model}

\subsection{Utilization Rate}

\begin{definition}[Utilization Rate]
For a lending pool with total deposits $S$ and total borrows $B$:
\begin{equation}
U = \frac{B}{S}
\end{equation}
\end{definition}

\subsection{Interest Rate Curve}

\begin{definition}[Kinked Interest Rate Model]
The borrow rate as a function of utilization:
\begin{equation}
r_b(U) = \begin{cases}
r_0 + U \cdot \frac{r_1 - r_0}{U^*} & \text{if } U \leq U^* \\
r_1 + (U - U^*) \cdot \frac{r_2 - r_1}{1 - U^*} & \text{if } U > U^*
\end{cases}
\end{equation}
where:
\begin{itemize}
\item $r_0$: Base rate (intercept)
\item $r_1$: Rate at optimal utilization
\item $r_2$: Maximum rate
\item $U^*$: Optimal utilization (kink point)
\end{itemize}
\end{definition}

\begin{figure}[h]
\centering
\begin{tabular}{|c|c|}
\hline
Parameter & Typical Value \\
\hline
$r_0$ & 0\% \\
$r_1$ & 4\% \\
$r_2$ & 100\% \\
$U^*$ & 80\% \\
\hline
\end{tabular}
\caption{Typical interest rate parameters}
\end{figure}

\subsection{Supply Rate}

\begin{theorem}[Supply Rate]
The supply rate is derived from the borrow rate:
\begin{equation}
r_s = r_b \cdot U \cdot (1 - \rho)
\end{equation}
where $\rho$ is the protocol reserve factor.
\end{theorem}

\begin{proof}
Total interest paid by borrowers: $r_b \cdot B$

Interest distributed to suppliers: $r_b \cdot B \cdot (1 - \rho)$

Per-unit supply rate: $\frac{r_b \cdot B \cdot (1 - \rho)}{S} = r_b \cdot U \cdot (1 - \rho)$
\end{proof}

\subsection{Compound Interest}

\begin{definition}[Interest Index]
The cumulative interest index at time $t$:
\begin{equation}
I(t) = I(t_0) \cdot \prod_{i=0}^{n-1} (1 + r_i \cdot \Delta t_i)
\end{equation}
For continuous compounding:
\begin{equation}
I(t) = I(t_0) \cdot \exp\left(\int_{t_0}^{t} r(\tau) \, d\tau\right)
\end{equation}
\end{definition}

\begin{proposition}[Debt Accumulation]
User debt at time $t$ given initial debt $D_0$ at time $t_0$:
\begin{equation}
D(t) = D_0 \cdot \frac{I(t)}{I(t_0)}
\end{equation}
\end{proposition}

\section{Risk Management}

\subsection{Oracle Design}

Price oracles are critical for accurate collateral valuation:

\begin{definition}[Time-Weighted Average Price]
TWAP over window $[t-w, t]$:
\begin{equation}
\text{TWAP}_w = \frac{1}{w} \int_{t-w}^{t} P(\tau) \, d\tau
\end{equation}
\end{definition}

\begin{proposition}[TWAP Manipulation Resistance]
The cost to manipulate TWAP by $\delta$ for window $w$:
\begin{equation}
\text{Cost} \geq \delta \cdot L \cdot w
\end{equation}
where $L$ is market liquidity. Longer windows increase manipulation cost.
\end{proposition}

\subsection{Supply and Borrow Caps}

\begin{definition}[Protocol Caps]
\begin{itemize}
\item Supply cap $S_{\max}$: Maximum total deposits
\item Borrow cap $B_{\max}$: Maximum total borrows
\end{itemize}
\end{definition}

These caps limit protocol exposure to any single asset.

\subsection{Reserve Factor}

\begin{proposition}[Reserve Accumulation]
Protocol reserves at time $t$:
\begin{equation}
R(t) = R(0) + \int_0^t \rho \cdot r_b(\tau) \cdot B(\tau) \, d\tau
\end{equation}
\end{proposition}

Reserves provide a buffer against bad debt and fund protocol development.

\section{Implementation on Lux}

\subsection{Storage Optimization}

Efficient on-chain storage is critical:

\begin{itemize}
\item \textbf{Packed Storage}: Multiple values in single 256-bit slot
\item \textbf{Index-Based Accounting}: Store indices, compute balances on-read
\item \textbf{Lazy Updates}: Batch interest accrual
\end{itemize}

\subsection{Core Borrow Function}

\begin{algorithm}
\caption{Borrow Execution}
\begin{algorithmic}[1]
\Require User $u$, asset $a$, amount $A$
\Ensure Loan disbursed or reverted
\State \textbf{accrue\_interest}$(a)$
\State $C \gets \text{get\_collateral\_value}(u)$
\State $D \gets \text{get\_debt\_value}(u)$
\State $D_{\text{new}} \gets D + A \cdot P_a$
\State \textbf{require} $D_{\text{new}} \leq C \cdot \lambda_{\text{weighted}}$
\State \textbf{require} $B_a + A \leq B_{\max,a}$
\State \textbf{transfer}$(a, u, A)$
\State \textbf{update\_debt}$(u, a, A)$
\end{algorithmic}
\end{algorithm}

\section{Governance and Parameter Updates}

\subsection{Parameter Sensitivity}

\begin{theorem}[Liquidation Threshold Impact]
Increasing $\theta$ by $\Delta\theta$ increases maximum borrowing capacity by:
\begin{equation}
\frac{\Delta B}{B} \approx \frac{\Delta\theta}{\theta}
\end{equation}
but also increases liquidation risk.
\end{theorem}

\subsection{Time-Locked Updates}

Critical parameters should have time-locked updates:
\begin{itemize}
\item Liquidation threshold: 7-day timelock
\item Interest rate model: 3-day timelock
\item Oracle sources: 7-day timelock
\end{itemize}

\section{Conclusion}

We have presented a complete mathematical framework for decentralized lending protocols. The framework ensures:

\begin{enumerate}
\item \textbf{Solvency}: Over-collateralization with asset-specific factors
\item \textbf{Prompt Liquidation}: Incentive-compatible liquidation mechanism
\item \textbf{Market-Driven Rates}: Kinked interest rate model balancing supply and demand
\end{enumerate}

The Lux blockchain's high throughput enables efficient liquidation processing and real-time interest accrual, making it well-suited for lending protocol deployment.

\begin{thebibliography}{9}
\bibitem{aave}
Aave Protocol Whitepaper. (2020). Aave V2.

\bibitem{compound}
Leshner, R., Hayes, G. (2019). Compound: The Money Market Protocol. Compound Whitepaper.

\bibitem{makerdao}
MakerDAO Team. (2017). The Dai Stablecoin System. MakerDAO Whitepaper.

\bibitem{euler}
Euler Finance. (2021). Euler: A permissionless lending protocol. Euler Whitepaper.

\bibitem{morpho}
Morpho Labs. (2022). Morpho: Peer-to-peer layer for lending pools. Morpho Whitepaper.
\end{thebibliography}

\end{document}
